\documentclass{beamer}

\usepackage[utf8]{inputenc}
\input{../helper}

\title{Introdução à Computação \\ estruturas básicas da linguagem C}
\author{Alexandre Rademaker}
\date{}

\begin{document}

\begin{frame}
 \maketitle
\end{frame}


\begin{frame}{estruturas de programação}

  \begin{itemize}
  \item sequência
  \item decisão
  \item repetição
  \end{itemize}

  O sequenciamento de comandos é definido pelo operador \ilcode{;},
  que termina cada instrução. Blocos de sequências são delimitados
  pelas \{ e \}.
  
\end{frame}

\begin{frame}[fragile,allowframebreaks]{condicionais e expressões lógicas}

  Escreva um programa para encontrar as raízes de uma equação do
  segundo grau:
  \[ a x^2 + b x + c = 0 \]

  Os coeficientes da equação são reais. O programa faz a alocação de 3
  posições de memória para esses coeficientes, inicializando-os com o
  valor zero.
  
  O programa efetua a leitura dos coeficientes através do teclado e,
  em seguida, calcula o valor das raízes existentes.
  
  Caso não existam raízes reais, o programa deve informar este fato ao
  usuário.

  \framebreak

  \lstinputlisting[style=tinyc,caption=src/raiz.c]{src/raiz.c}

  \framebreak

  \lstinputlisting[style=tinyc,caption=src/parity0.c]{src/parity0.c}

  O operador \ilcode{\%} nos dá o 'resto' da divisão. Vamos compilar e
  executar?

  \framebreak
  
  \lstinputlisting[style=tinyc,caption=src/agree.c]{src/agree.c}

  Para caracteres usamos `single quotes' e para strings `double
  quotes'.

  \framebreak

  \begin{itemize}
  \item operadores relacionais: \ilcode{a > b}, \ilcode{a < b},
    \ilcode{a >= b}, \ilcode{a <= b}, \ilcode{a == b} e \ilcode{a !=
      b} (diferente).

  \item operadores lógicos: \ilcode{!p} (not), \ilcode{p \|\| q} (or),
    \ilcode{p \&\& q} (and). Onde \ilcode{p} e \ilcode{q} são
    expressões boleanas.
  \end{itemize}

  Em C os inteiros também são usados como valores booleanos: qualquer
  valor não nulo (1) representa \textbf{true} e 0 representa
  \textbf{false}. Biblioteca \ilcode{stdbool.h} declara estas
  constantes.
  
  
  \framebreak

  \lstinputlisting[style=tinyc,caption=src/area.c]{src/area.c}

\end{frame}


\begin{frame}[allowframebreaks,fragile]{Repetição, repetição \ldots}

  Todo bom programador é `preguiçoso'! Péssimo estilo de código. \vspace{.5cm}

  \lstinputlisting[style=tinyc,caption=src/meow0.c]{src/meow0.c}

  \framebreak

  \lstinputlisting[style=tinyc,caption=src/meow1.c]{src/meow1.c}
    
  \lstinputlisting[style=tinyc,caption=src/meow2.c]{src/meow2.c}

  \framebreak

  Podemos repetir para sempre \ldots

  \begin{lstlisting}[style=normalc]
    while (true)
    {
      printf("meow\n");
    }
  \end{lstlisting}

  Precisamos de algumas bibliotecas adicionais para poder usar o valor
  \ilcode{true}, a \ilcode{stdbool.h}.

  Use \ilcode{CTRL+C} para terminar!

  \framebreak

  \lstinputlisting[style=tinyc,caption=src/guess.c]{src/guess.c}
  
\end{frame}


\end{document}

%%% Local Variables:
%%% mode: latex
%%% TeX-master: t
%%% End:
